% Fichier complété par tools/complete_missing_corrections.py.
% Source : DYN/DYN-01/63_BancHydraulique/63_BancHydraulique.tex
% Statut : rédigé (7 question(s)).
\begin{corrigebox}[Corrigé rédigé]
\CorrigeQuestion{1}
Pour un profil trapézoïdal, la distance rapide est l'aire sous la
                courbe de vitesse : $c_{\rm rap}=V(t_r+t_a)$ et la distance lente
                $c_{\rm lent}=V_l(t_l+t_a)$, avec demi-aires $Vt_a/2$ à chaque rampe.
                Les expressions exactes s'obtiennent en additionnant les trapèzes du
                chronogramme.
\CorrigeQuestion{2}
On résout les deux équations de course avec les durées imposées.
                L'accélération du chariot est ensuite $a=V/t_a$ pendant la montée et
                $-a$ pendant la descente.
\CorrigeQuestion{3}
Si le moteur entraîne une vis de pas $p$ par un réducteur de rapport
                $r=\omega_{\rm vis}/\omega_m$, alors
                $V=p\omega_{\rm vis}/(2\pi)=pr\omega_m/(2\pi)$, donc
                $\omega_m=2\pi V/(pr)$.
\CorrigeQuestion{4}
$\omega_{m,\max}=2\pi V_{\max}/(pr)$ et
                $\dot\omega_m=2\pi a/(pr)$.
\CorrigeQuestion{5}
$E_c=\tfrac12MV^2+\tfrac12I_m\omega_m^2+
                \tfrac12I_r\omega_r^2$ (on ajoute de même chaque pièce tournante).
\CorrigeQuestion{6}
En écrivant $E_c=\tfrac12J_{\rm eq}\omega_m^2$ :
                $J_{\rm eq}=I_m+I_r(\omega_r/\omega_m)^2+M(V/\omega_m)^2$.
\CorrigeQuestion{7}
Le théorème de l'énergie cinétique donne
                $C_m\omega_m-FV=J_{\rm eq}\omega_m\dot\omega_m$, donc
                $\boxed{C_m=J_{\rm eq}\dot\omega_m+F(V/\omega_m)}$.
\end{corrigebox}
